\documentclass[12pt,letterpaper]{screenplay}

% Title page information
\title{Your Screenplay Title}
\author{Your Name}
\address{Your Street Address \\ City, State ZIP \\ Phone \\ email@example.com}

\begin{document}

% Use \coverpage for full title page with contact info
% Or use \nicholl for simple title-only page
\coverpage

\fadein

\intslug[day]{Location name}

Action description of what happens in the scene. Describe the setting, character actions, and visual details here.

\begin{dialogue}{CHARACTER NAME}
What the character says as dialogue.
\end{dialogue}

\begin{dialogue}[emotional direction]{CHARACTER NAME}
Dialogue with an optional direction that appears in parentheses below the character name.
\end{dialogue}

Action description of the next beat in the scene.

\begin{dialogue}{CHARACTER NAME}
First part of dialogue.
\paren{reaction or action}
Continuing dialogue after the parenthetical direction.
\end{dialogue}

\extslug[night]{Exterior location name}

Action description for the exterior scene.

\begin{dialogue}{CHARACTER ONE}
First character's dialogue.
\end{dialogue}

\begin{dialogue}{CHARACTER TWO}
Second character's response.
\end{dialogue}

% Use \centretitle for flashbacks, time jumps, or title cards
\centretitle{Three Years Earlier}

\intslug[day]{Flashback location}

Action description in the flashback scene.

\begin{dialogue}{CHARACTER NAME}
Dialogue in the flashback.
\end{dialogue}

\extslug[night]{Another location}

Action description. Use \pov\ for point-of-view shots.

\pov\ Character's perspective of what they see.

\revert\ Back to normal angle.

More action description continuing the scene.

% \intercut produces "INTERCUT WITH:" flush right
\intercut

Action description of intercutting between scenes or locations.

\fadeout

\theend

\end{document}
